\section{Procesamiento distribuido}
\label{sec:procesamiento}

\subsection{Metodología}

Se definió un único conjunto de 12 consultas de negocio, y se implementó exactamente el mismo conjunto en los cuatro frameworks pedidos por el enunciado: Polars, Dask, Modin y Apache Spark. La razón de usar el mismo set de consultas en los cuatro, en vez de un set distinto por framework, es que el objetivo de esta sección es comparar cómo cada herramienta resuelve el mismo problema, no evaluar cuatro análisis de negocio independientes. Cada framework lee el CSV original desde cero y corre el pipeline completo de forma independiente (no reutiliza resultados intermedios de otro framework).

\begin{table}[H]
\centering
\small
\begin{tabular}{clp{7.5cm}}
\toprule
\textbf{\#} & \textbf{Categoría(s)} & \textbf{Consulta} \\
\midrule
1 & Limpieza & Eliminar filas con \texttt{Description} nula o \texttt{Price} $\leq$ 0 \\
2 & Duplicados & Detectar y eliminar filas exactamente duplicadas \\
3 & Transformación & Crear \texttt{TotalPrice}, parsear fecha, derivar \texttt{Year}/\texttt{Month}, \texttt{IsCancelled} \\
4 & Filtrado & Quedarse solo con ventas válidas (cantidad y precio positivos, no canceladas) \\
5 & Agrupación + Agregación & Top 10 productos por revenue total \\
6 & Agrupación + Agregación & Top 10 productos por unidades vendidas \\
7 & Agrupación + Ordenamiento & Top 10 clientes por gasto total \\
8 & Agrupación + Ordenamiento & Revenue y ticket promedio por país (top 10) \\
9 & Métrica temporal & Revenue mensual y variación \% mes a mes \\
10 & Métrica + Filtrado & Tasa de cancelación por país (top 10) \\
11 & Agrupación + Métrica & Segmentación de clientes en 3 tiers de gasto \\
12 & Métrica & Ticket promedio (AOV) global y su evolución mensual \\
\bottomrule
\end{tabular}
\caption{Las 12 consultas implementadas en cada uno de los 4 frameworks}
\label{tab:consultas}
\end{table}

Las consultas 1 a 4 preparan el dataset (limpieza, deduplicación, transformación, filtrado), las consultas 5 a 12 son las preguntas de negocio propiamente dichas, y se calculan sobre el subconjunto ya limpio de ventas válidas (salvo la consulta 10, que necesita ver también las transacciones canceladas para calcular una tasa).

\subsection{Polars}

Polars se implementó usando su API de expresiones (\texttt{with\_columns}, \texttt{filter}, \texttt{group\_by().agg()}), en modo eager. A diferencia de los otros tres frameworks, Polars no tiene ningún modo distribuido multi-máquina en su versión de código abierto: es una librería de una sola máquina por diseño, que paraleliza internamente entre los núcleos del procesador mediante su motor en Rust. Esto no es una limitación del proyecto: para consultas de este tamaño (~1 millón de filas), Polars resultó ser, en la práctica, el más simple de implementar y uno de los más rápidos localmente.

Se usó \texttt{qcut} para la segmentación de clientes (consulta 11) y \texttt{pct\_change} para la variación mensual (consulta 9), ambos disponibles de forma nativa en la API de Polars.

Como referencia, así se ve la consulta 5 (top 10 productos por revenue) en Polars:

\begin{verbatim}
top10 = (
    df.group_by("Description")
    .agg(pl.col("TotalPrice").sum())
    .sort("TotalPrice", descending=True)
    .head(10)
)
\end{verbatim}

\subsection{Apache Spark}

Spark se implementó con la API de DataFrame (\texttt{groupBy().agg()}, funciones del módulo \url{pyspark.sql.functions}), usando una función de ventana (\texttt{Window} + \texttt{lag}) para la variación mensual y \texttt{approxQuantile} para la segmentación de clientes, ya que Spark no tiene un equivalente directo a \texttt{qcut}.

Spark es, de los cuatro, el único que se ejecutó también sobre el clúster real de Dataproc en modo genuinamente distribuido (ver Sección \ref{sec:gcp-cluster-run}): al correrlo con \texttt{spark-submit} sin forzar un \texttt{master} local, el job se somete a YARN y se reparte entre el nodo master y los dos workers.

Como referencia, así se ve la consulta 5 (top 10 productos por revenue) en Spark:

\begin{verbatim}
top10 = (
    df.groupBy("Description")
    .agg(F.sum("TotalPrice").alias("TotalPrice"))
    .orderBy(F.desc("TotalPrice"))
    .limit(10)
)
\end{verbatim}

\subsection{Modin}

Modin se implementó usando \texttt{modin.pandas}, que expone una API prácticamente idéntica a pandas, de hecho, la implementación es casi un calco de la de Dask, quitando las llamadas a \texttt{.compute()}. Por debajo, Modin distribuye el trabajo usando Ray como motor de ejecución. Al importar \texttt{modin.pandas} sin conectarse explícitamente a un clúster Ray existente, Ray arranca automáticamente una instancia local sobre la máquina donde corre el script.

Como referencia, así se ve la consulta 5 (top 10 productos por revenue) en Modin, prácticamente idéntica a pandas:

\begin{verbatim}
resultado = df.groupby("Description")["TotalPrice"].sum()
top10 = resultado.nlargest(10)
\end{verbatim}

\subsection{Dask}

Dask se implementó con \texttt{dask.dataframe}, cuya API está deliberadamente inspirada en pandas. Un detalle técnico relevante: al leer el CSV con un \texttt{blocksize} demasiado grande, Dask particionaba el archivo en una sola partición, lo que anulaba cualquier paralelismo real, se ajustó el \texttt{blocksize} para obtener varias particiones reales. También se agregó \texttt{.persist()} después de las etapas de limpieza y filtrado, para que el pipeline pesado no se recomputara desde cero en cada una de las ocho consultas de negocio restantes (Dask es perezoso por defecto).

Como referencia, así se ve la consulta 5 (top 10 productos por revenue) en Dask, con el \texttt{.compute()} explícito que dispara la ejecución:

\begin{verbatim}
resultado = df.groupby("Description")["TotalPrice"].sum().compute()
top10 = resultado.nlargest(10)
\end{verbatim}

\subsection{Resultados y comparación}

Las cuatro implementaciones, corriendo de forma completamente independiente, producen resultados consistentes entre sí en las 12 consultas: mismo top de productos y clientes, mismo ticket promedio, mismas tasas de cancelación. La única excepción es la consulta 11 (segmentación de clientes), donde Polars ubica a un cliente en un tier distinto al de los otros tres frameworks, el detalle de esa diferencia puntual se explica en la Sección \ref{sec:operaciones}. Fuera de ese caso, las cuatro coinciden exactamente, lo que confirma que resuelven correctamente el mismo problema, solo con sintaxis y motor de ejecución distintos.

La Tabla \ref{tab:cobertura} muestra que las 12 consultas de la Tabla \ref{tab:consultas} fueron implementadas y ejecutadas con éxito en los 4 frameworks, con resultado numérico coincidente en 11 de las 12 (la excepción es la consulta 11, ver Sección \ref{sec:operaciones}).

\begin{table}[H]
\centering
\small
\begin{tabular}{clcccc}
\toprule
\textbf{\#} & \textbf{Consulta} & \textbf{Polars} & \textbf{Dask} & \textbf{Modin} & \textbf{Spark} \\
\midrule
1 & Limpieza & $\checkmark$ & $\checkmark$ & $\checkmark$ & $\checkmark$ \\
2 & Duplicados & $\checkmark$ & $\checkmark$ & $\checkmark$ & $\checkmark$ \\
3 & Transformación & $\checkmark$ & $\checkmark$ & $\checkmark$ & $\checkmark$ \\
4 & Filtrado & $\checkmark$ & $\checkmark$ & $\checkmark$ & $\checkmark$ \\
5 & Top 10 productos por revenue & $\checkmark$ & $\checkmark$ & $\checkmark$ & $\checkmark$ \\
6 & Top 10 productos por unidades & $\checkmark$ & $\checkmark$ & $\checkmark$ & $\checkmark$ \\
7 & Top 10 clientes por gasto & $\checkmark$ & $\checkmark$ & $\checkmark$ & $\checkmark$ \\
8 & Revenue y ticket por país & $\checkmark$ & $\checkmark$ & $\checkmark$ & $\checkmark$ \\
9 & Revenue mensual y variación & $\checkmark$ & $\checkmark$ & $\checkmark$ & $\checkmark$ \\
10 & Tasa de cancelación por país & $\checkmark$ & $\checkmark$ & $\checkmark$ & $\checkmark$ \\
11 & Segmentación de clientes & $\checkmark^{\ast}$ & $\checkmark$ & $\checkmark$ & $\checkmark$ \\
12 & Ticket promedio (AOV) & $\checkmark$ & $\checkmark$ & $\checkmark$ & $\checkmark$ \\
\bottomrule
\end{tabular}
\caption{Cobertura de las 12 consultas por framework, implementadas y ejecutadas con éxito en los 4, con resultado numérico coincidente en 11 de las 12 filas}
\label{tab:cobertura}
\end{table}

{\footnotesize $\ast$Polars ubica a un cliente en un tier de gasto distinto al de los otros tres en la consulta 11, ver el detalle en la Sección \ref{sec:operaciones}.}

\subsubsection*{Corrida sobre el clúster real de Dataproc}
\label{sec:gcp-cluster-run}

Además de validarse en local, los cuatro frameworks se corrieron también sobre el clúster real de Dataproc descrito en la Sección \ref{sec:gcp}, leyendo el dataset directamente desde el bucket. Los cuatro reprodujeron exactamente los mismos resultados que en local. El output completo de cada corrida está disponible en \texttt{docs/evidencia/} (un archivo \texttt{-cluster-output.txt} por framework).

Vale la pena aclarar que, de los cuatro, solo Spark aprovecha genuinamente los tres nodos del clúster: Dataproc ya trae YARN configurado de fábrica, así que basta con no forzar un \texttt{master} local en el código. Dask y Modin, al no conectarse explícitamente a un scheduler distribuido (Dask) o a un clúster Ray (Modin) montado entre los nodos, corren en el nodo master usando sus propios núcleos, siguen ejecutándose sobre infraestructura de Dataproc, pero sin repartirse entre los tres nodos. Montar un clúster Dask o Ray real entre los tres nodos es posible, pero es trabajo de infraestructura adicional que no aporta valor proporcional al alcance de este proyecto.

\subsection{Comparación de rendimiento}
\label{sec:benchmark}

El objetivo del proyecto es correr el procesamiento distribuido sobre GCP, así que la comparación de rendimiento se mide sobre el clúster real de Dataproc descrito en la Sección \ref{sec:gcp}. Se corrió el mismo pipeline completo (carga del CSV más las 12 consultas) una vez por framework, sobre el mismo clúster de prueba: Polars, Dask y Modin se ejecutaron directamente en el nodo master (\texttt{n4-standard-2}, 2 vCPU, 8 GB RAM), y Spark se envió como job real de YARN (\texttt{gcloud dataproc jobs submit pyspark}), repartido entre el master y los 2 workers. Se midió el tiempo de pared de cada etapa con \texttt{time.perf\_counter}. El script usado es \url{processing/run_benchmark.py} (y su variante \url{processing/spark/benchmark_cluster.py} para el envío a YARN), y la salida cruda de cada corrida está en \texttt{docs/evidencia/benchmark/*-cluster.json}.

\begin{table}[H]
\centering
\begin{tabular}{lrl}
\toprule
\textbf{Framework} & \textbf{Tiempo total (s)} & \textbf{Dónde corrió} \\
\midrule
Polars & 1.76 & nodo master (proceso local) \\
Dask & 11.25 & nodo master (proceso local) \\
Modin & 12.40 & nodo master (proceso local) \\
Spark & 170.65 & job de YARN, 3 nodos \\
\bottomrule
\end{tabular}
\caption{Tiempo total del pipeline completo, corrido en el clúster real de Dataproc}
\label{tab:benchmark-cluster}
\end{table}

El resultado es honesto pero no es el esperado: Spark corriendo genuinamente distribuido en YARN (170.65 s) resultó el más lento de los cuatro, no el más rápido. Esto no contradice que Spark sea el único que reparte trabajo entre nodos (ver Sección \ref{sec:gcp-cluster-run}), confirma algo distinto: para un dataset de $\sim$1 millón de filas, en un clúster de prueba con máquinas modestas (\texttt{n4-standard-2}, 2 vCPU cada una), el overhead de coordinar tareas y mover datos entre 3 nodos por red pesa más que el beneficio de paralelizar el trabajo. Ese overhead se amortiza con datasets mucho más grandes o clústers con más núcleos por nodo, no con este volumen de datos ni este hardware de prueba.

La Tabla \ref{tab:benchmark-etapas-cluster} desglosa el tiempo por etapa en el clúster.

\begin{table}[H]
\centering
\small
\begin{tabular}{lrrrr}
\toprule
\textbf{Etapa} & \textbf{Polars} & \textbf{Dask} & \textbf{Modin} & \textbf{Spark} \\
\midrule
Inicialización de motor & --- & --- & --- & 8.87 \\
Carga del CSV & 0.41 & 0.02 & 5.15 & 2.86 \\
1. Limpieza & 0.13 & 1.93 & 0.92 & 9.03 \\
2. Duplicados & 0.28 & 3.18 & 1.17 & 6.89 \\
3. Transformación & 0.35 & 4.35 & 1.07 & 5.38 \\
4. Filtrado & 0.03 & 0.87 & 0.76 & 32.45 \\
5. Top productos (revenue) & 0.10 & 0.06 & 0.34 & 20.34 \\
6. Top productos (unidades) & 0.06 & 0.04 & 0.30 & 12.11 \\
7. Top clientes & 0.05 & 0.05 & 0.45 & 11.17 \\
8. Revenue por país & 0.09 & 0.27 & 0.59 & 12.64 \\
9. Revenue mensual & 0.05 & 0.04 & 0.13 & 6.62 \\
10. Tasa de cancelación & 0.04 & 0.06 & 0.48 & 6.38 \\
11. Segmentación de clientes & 0.07 & 0.04 & 0.46 & 19.85 \\
12. Ticket promedio (AOV) & 0.08 & 0.34 & 0.58 & 24.93 \\
\bottomrule
\end{tabular}
\caption{Tiempo por etapa del pipeline, en segundos, corrido en el clúster real de Dataproc}
\label{tab:benchmark-etapas-cluster}
\end{table}

La etapa 4 (Filtrado) es la más reveladora: en Spark salta a 32.45 s, muy por encima de las demás etapas, porque ahí es donde el código llama a \texttt{.cache()} sobre \texttt{validas} por primera vez, lo que fuerza a materializar y mover datos entre los 2 executors distribuidos en los workers, el costo real de la distribución aparece exactamente ahí, en el shuffle de red entre nodos. Las consultas 5, 6, 7, 8, 11 y 12 también son caras en Spark porque cada una dispara una acción distribuida nueva sobre ese DataFrame ya cacheado. Polars, Dask y Modin, al no repartirse entre nodos, quedan limitados a los 2 vCPU / 8 GB del nodo master, muy por debajo de lo que tendrían disponible si se distribuyeran entre los 3 nodos del clúster.
